% Standalone supplementary document — open in the file tree and click Compile.
\documentclass[11pt,a4paper]{article}

\usepackage[utf8]{inputenc}
\usepackage[T1]{fontenc}
\usepackage[margin=1in]{geometry}
\usepackage{microtype}
\usepackage{amsmath,amssymb}
\usepackage{booktabs}
\usepackage{siunitx}
\usepackage{hyperref}

\title{\textbf{Supplementary Information}\\[4pt]
       \large Reproducible Multi-File \LaTeX{} Workflows}
\author{First Author \and Second Author}
\date{\today}

\begin{document}
\maketitle

\section*{Overview}
This PDF is built from \texttt{supplement.tex} alone.
The main manuscript lives in \texttt{main.tex}; conference slides in \texttt{slides.tex}.
All three share \texttt{references.bib}.

\section{Extended experimental results}
Table~\ref{tab:extended} reports per-seed metrics omitted from the main paper for brevity.

\begin{table}[htbp]
  \centering
  \caption{Extended benchmark (five seeds).}
  \label{tab:extended}
  \begin{tabular}{@{}lccc@{}}
    \toprule
    \textbf{Seed} & \textbf{Accuracy} & \textbf{F1} & \textbf{Time (s)} \\
    \midrule
    1 & 0.908 & 0.891 & 44 \\
    2 & 0.917 & 0.902 & 41 \\
    3 & 0.911 & 0.895 & 46 \\
    4 & 0.919 & 0.904 & 42 \\
    5 & 0.914 & 0.898 & 43 \\
    \bottomrule
  \end{tabular}
\end{table}

\section{Additional related work}
We include further citations for reproducibility tooling~%
\cite{example2021survey,example2025toolkit} and publishing practice~%
\cite{example2022conference}.

\section{Proof sketch}
For completeness, the regularised objective from the main text satisfies
$\nabla_\theta \mathcal{L} = \frac{2}{N} X^\top (X\theta - y) + 2\lambda\theta$ under standard assumptions.

\bibliographystyle{IEEEtran}
\bibliography{references}

\end{document}
